% !TEX TS-program = LuaLaTeX
% !TEX encoding = UTF-8 Unicode
\documentclass[%
   a4paper,
   12pt,
   twoside,
   tipotesi=scudo,
   ]{toptesi}
 
%%%% Use the following package if and only if you want to produce
%%%% an archivable document according to standard PDF/A-1b
%%%% No need to load package imakeidx, because it has already been
%%%% loaded by the specific module toptesi-scudo.
\usepackage[a-1b]{pdfx}
%%%%%% Read the English documentation of TOPtesi in order to check
%%%%%% the special attention needed to produce ISO compliant
%%%%%% archivable files

%%%%

\ifPDFTeX
    \usepackage[utf8]{inputenc}% 
    \usepackage[T1]{fontenc}
\fi

\errorcontextlines=9% more information on the console in case of errors

%%%% Specify fonts here; chose one among these fonts by leaving just
%%%% one line without initial comment character.
%%%% With LuaLaTeX or XeLaTeX don't change fonts
    \usepackage{fontspec}
    \defaultfontfeatures{Ligatures=TeX}
    \setmainfont{Libertinus Serif}
    \setsansfont{Libertinus Sans}
    \setmonofont{Libertinus Mono}[Scale=MatchLowercase]
    \usepackage{unicode-math}% add special math stile option here
                             % for example [style=ISO]
% chose one math font by commenting the line of the undesired font
    \setmathfont{XITS Math}




\raggedbottom
\begin{document}


% The contents of this ThesisTitlePage environment may be written
% in a configuration file named exactly the same as the thesis main
% file, but with extension .cfg. If similar commands with different
% data are written within this environment, such data prevail on
% those read from the configuration file.
\begin{ThesisTitlePage}
%\renewcommand\fivestars{\rule{0pt}{\baselineskip}\rule{30mm}{0.4pt}}
% Use the optional command to set a different School logo
\PhDschoolLogo{TITDocScCropped}
% Doctorate course name; mandatory
\ProgramName{Energy Enginering}
% Cycle ordinal number; optional.
% You can write 29.th, or 29\ap{th}, or 29\textsuperscript{th}, or ...
\CycleNumber{29.th}
% PhD candidate name; mandatory
\author{Mario Rossi}
% Dissertation title; mandatory
\title{Writing your Doctoral~Thesis\\with \LaTeX}
% Dissertation subtitle: optional.
% It might be useful only if the actual full title is too long.
\subtitle{Typeset with the TOPtesi class}
% The supervisor(s) label; optional; default value "Supervisor:".
% You can change it to plural as in this example, where the colon has
% been eliminated.
%\NSupervisor{Supervisor}{Supervisors}
%
% The SupervisorNumber may contain a value such as 0, 1, or any
% number higher than 1. If 0 is specified, no label is typeset
% over the supervisor(s) list; if 1 is specified then the singular
% form is used: if any value higher than 1 is specified, the plural
% form is used.
\SupervisorNumber{2}
% List of supervisors with academic title, name(s), surname(s),
% and function; mandatory
\SupervisorList{%
    Prof.~A.B., Supervisor\\
    Prof.~C.D. Co-supervisor}
% Name of the examining committee: optional. 
% Default value "Doctoral Examination Committee"
%\Nexaminationcommittee{Doctoral examination committee}
% List of the  examining committee members: mandatory if the above label
% is not empty.
\ExaminerList{%
Prof.~A.B., Referee, University of \dots\\
Prof.~C.D., Referee, University of \dots\\
Prof.~E.F., University of \dots\\
Prof.~G.H., University of \dots\\
Prof.~I.J., University of \dots}
% Name of the institution where the examination takes place; optional.
% Default value: "Politecnico di Torino"
%\Nlocation{Politecnico di Torino}
% Examination date: mandatory, although the way to write it is free.
\ExaminationDate{February 29, 2123}
% Disclaimer with signature; optional. Default text as
% in the following lines. 
%\Signature{%
%\begin{flushright}
%\parbox{0.5\textwidth}{\centering
%\dotfill\\
%Mario Rossi\\
%Turin, February 29, 2123
%}
%\end{flushright}}
\end{ThesisTitlePage}
\end{document}